\documentclass[12pt,a4paper]{article}
\usepackage{cmap} % Makes the PDF copiable. See http://tex.stackexchange.com/a/64198/25761
\usepackage[T1]{fontenc}
\usepackage[brazil]{babel}
\usepackage[utf8]{inputenc}
\usepackage{amsmath}
\usepackage{amsfonts}
\usepackage{amssymb}
\usepackage{amsthm}
\usepackage{textcomp} % \degree
\usepackage{gensymb} % \degree
\usepackage[usenames,svgnames,dvipsnames]{xcolor}
\usepackage{hyperref}
\usepackage{graphicx}
\usepackage[margin=2cm]{geometry}

\hypersetup{
    colorlinks = true,
    allcolors = {blue}
}

\newcommand*\diff{\mathop{}\!\mathrm{d}}
\newcommand*\sen{\operatorname{sen}}
\newcommand*\senh{\operatorname{senh}}
\newcommand*\cotg{\operatorname{cotg}}
\newcommand*\arctg{\operatorname{arctg}}

\newcommand*\prova{Prova IV}
\newcommand*\turma{TADS121-01C}
\newcommand*\disciplina{CDI0001}
\newcommand*\eu{Helder G. G. de Lima}
\newcommand*\data{05/12/2016}

\author{\eu}
\title{\prova - \disciplina}
\date{\data}

\begin{document}
\thispagestyle{empty}
\newgeometry{margin=2cm,bottom=0.5cm}
\begin{center}
\includegraphics[width=9.0cm]{marca}
\noindent\begin{tabular}{l c c r}
  \textbf{\disciplina (\turma)}
& \textbf{\prova}
& \textbf{\data}
\end{tabular}
\\ Prof. \eu\footnote{
Este é um material de acesso livre distribuído sob os termos da licença \href{https://creativecommons.org/licenses/by-sa/4.0/deed.pt_BR}{Creative Commons BY-SA 4.0}.}
\end{center}

\noindent Nome do(a) aluno(a): \rule{13cm}{0.01cm}

%\section*{Instruções}
\begin{center}\fbox{
\begin{minipage}{14cm}

{\footnotesize
\begin{itemize}
\renewcommand{\theenumi}{\Roman{enumi}}
\item Identifique-se em todas as folhas.
\item Mantenha o celular e os demais equipamentos eletrônicos desligados durante a prova.
\item Justifique cada resposta com cálculos ou argumentos baseados na teoria estudada.
\end{itemize}
}

\end{minipage}
}
\end{center}


%\section*{Questões}

\begin{enumerate}
\item Calcule apenas 5 das 6 integrais a seguir, deixando explícito como obteve todas as integrais que não forem imediatas:
\begin{enumerate}
\item (2,0) $\displaystyle \int \frac{12 \cos(3x)}{\sec(15x)} \diff x$
\item (2,0) $\displaystyle \int e^{2x} \cosh(3x) \diff x$
\item (2,0) $\displaystyle \int 3x \sqrt{5x^2+8} -\cotg(x) \diff x$
\item (2,0) $\displaystyle \int \frac{25}{(x^2+9)(x^2+4)} \diff x$
\item (2,0) $\displaystyle \int \frac{e^{(x+3)}}{1+e^{(2x)}} \diff x$
\item (2,0) $\displaystyle \int \frac{9 x^2 - 14 x + 2}{x^3-2x^2+x} \diff x$
\end{enumerate}
\end{enumerate}

\begin{center}
BOA PROVA!!!
\end{center}

\newpage
\restoregeometry
\section*{Respostas e observações}
\begin{enumerate}
\item
\begin{enumerate}
\item \textbf{Solução 1}:
\begin{align*}
    \int \frac{12 \cos(3x)}{\sec(15x)} \diff x
& = \int 12 \cos(3x) \cos(15x) \diff x
  = \int \frac{12}{2}\left[ \cos(3x-15x) + \cos(3x+15x) \right] \diff x\\
& = 6 \int \cos(-12x) + \cos(18x) \diff x
  = 6 \int \cos(12x) \diff x + 6\int \cos(18x) \diff x\\
& = 6 \frac{ \sen(12x) }{12} + 6\frac{ \sen(18x) }{18} + c
  = \frac{ \sen(12x) }{2} + \frac{ \sen(18x) }{3} + c.
\end{align*}

\textbf{Solução 2}:
Seja $u=\cos(3x)$,
     $\diff u=-3\sen(3x) \diff x$,
     $\diff v=\cos(15x) \diff x$ e
     $v=\dfrac{\sen(15x)}{15}$. Então:
\begin{align*}
    \int \frac{12 \cos(3x)}{\sec(15x)} \diff x
& = 12 \int \cos(3x) \cos(15x) \diff x \\
& = 12 \left[ \cos(3x)\frac{\sen(15x)}{15}
  - \int \frac{\sen(15x)}{15} (-3\sen(3x)) \diff x \right]\\
& = \frac{4}{5}\cos(3x)\sen(15x)
  + \frac{12}{5}\int \sen(15x) \sen(3x) \diff x.
\end{align*}
Sendo $U=\sen(3x)$,
     $\diff U=3\cos(3x) \diff x$,
     $\diff V=\sen(15x) \diff x$,
     $V=-\dfrac{\cos(15x)}{15}$, tem-se:
\begin{align*}
    \frac{12}{5}\int \sen(15x) \sen(3x) \diff x
& = \frac{12}{5} \left[ \sen(3x)\left(-\dfrac{\cos(15x)}{15}\right)
                 - \int - \dfrac{\cos(15x)}{15} 3\cos(3x) \diff x \right]\\
& = -\frac{4}{25} \sen(3x)\cos(15x) + \frac{12}{25}\int \cos(3x)\cos(15x) \diff x.
\end{align*}
Assim,
\begin{align*}
12 \int \cos(3x) \cos(15x) \diff x
& = \frac{4}{5}\cos(3x)\sen(15x) - \frac{4}{25} \sen(3x)\cos(15x)\\
& +\frac{1}{25} \left( 12\int \cos(3x)\cos(15x) \right),
\end{align*}
e os termos envolvendo a integral procurada podem ser agrupados no primeiro membro, levando em conta que $1-\frac{1}{25} = \frac{24}{25}$, para obter
\[
\frac{24}{25} \cdot 12 \int \cos(3x) \cos(15x) \diff x
=
\frac{4}{5}\cos(3x)\sen(15x) - \frac{4}{25} \sen(3x)\cos(15x) + c_1.
\]
Portanto,
\[
12 \int \cos(3x) \cos(15x) \diff x
=
\frac{5}{6}\cos(3x)\sen(15x) - \frac{1}{6} \sen(3x)\cos(15x) + c_2.
\]
\textbf{Observação}: Esta função é a mesma obtida na solução 1, como pode ser verificado por meio da identidade trigonométrica $\sen(a)\cos(b)=\frac{1}{2}\left[ \sen(a-b) + \sen(a+b)\right]$.
\item \textbf{Solução 1}:
\begin{align*}
\int e^{2x} \cosh(3x) \diff x
& = \int e^{2x} \left( \frac{e^{3x}+e^{-3x}}{2} \right)\diff x
  = \int \frac{e^{5x} + e^{-x}}{2} \diff x\\
& = \frac{1}{2}\int e^{5x} \diff x + \frac{1}{2}\int e^{-x} \diff x
  = \frac{1}{10}\int e^{5x} \cdot 5\diff x
  - \frac{1}{2}\int e^{-x} (-1)\diff x\\
& = \frac{1}{10}\int e^u \diff u
  - \frac{1}{2}\int e^{v} \diff v
  = \frac{e^u}{10} - \frac{e^v}{2}  + c
  = \frac{e^{5x}}{10} - \frac{e^{-x}}{2} + c.
\end{align*}
\textbf{Solução 2}: Tomando
$u=e^{2x}$,
$\diff u=2 e^{2x} \diff x$,
$\diff v=\cosh(3x)\diff x$ e
$v=\frac{\senh(3x)}{3}$, tem-se:
\begin{align*}
\int e^{2x} \cosh(3x) \diff x
& = e^{2x} \frac{\senh(3x)}{3} - \int \frac{\senh(3x)}{3} 2 e^{2x} \diff x\\
& = \frac{1}{3}e^{2x}\senh(3x) - \frac{2}{3}\int e^{2x} \senh(3x) \diff x.
\end{align*}
Agora, escolhendo
$U=e^{2x}$,
$\diff U=2 e^{2x} \diff x$,
$\diff V=\senh(3x)\diff x$ e
$V=\frac{\cosh(3x)}{3}$, resulta:
\begin{align*}
\\
\int e^{2x} \cosh(3x) \diff x
& = \frac{1}{3}e^{2x}\senh(3x) - \frac{2}{3}\left[
e^{2x}\frac{\cosh(3x)}{3}
-
\int \frac{\cosh(3x)}{3} 2 e^{2x} \diff x
\right]\\
& = \frac{1}{3}e^{2x}\senh(3x)
- \frac{2}{9} e^{2x}\cosh(3x)
+ \frac{4}{9} \int e^{2x} \cosh(3x) \diff x.
\end{align*}
Agrupando os termos que envolvem a integral procurada, tem-se:
\begin{align*}
\frac{5}{9} \int e^{2x} \cosh(3x) \diff x
& = \frac{1}{3}e^{2x}\senh(3x) - \frac{2}{9} e^{2x}\cosh(3x) + c_1.
\end{align*}
Portanto,
\begin{align*}
\int e^{2x} \cosh(3x) \diff x
= \frac{3}{5} e^{2x}\senh(3x) - \frac{2}{5} e^{2x}\cosh(3x)
= \frac{e^{5x}}{10} + \frac{e^{-x}}{2} + c_2.
\end{align*}

\textbf{Observação}: Esta função é a mesma obtida na solução 1, como pode ser verificado por meio da definição do cosseno hiperbólico.

\item Tem-se $
\int 3x \sqrt{5x^2+8} -\cotg(x) \diff x
= \int 3x \sqrt{5x^2+8} \diff x
 -\int \cotg(x) \diff x
$. A primeira integral pode ser resolvida utilizando a substituição $t = 5x^2+8$, $\diff t = 10x \diff t$:
\begin{align*}
\int 3x \sqrt{5x^2+8} \diff x
& = \frac{3}{10} \int \sqrt{5x^2+8} \cdot 10x \diff x
= \frac{3}{10} \int \sqrt{t} \diff t\\
&= \frac{3}{10} \left( \frac{2}{3}\sqrt{t^3}\right) + k_1
= \frac{1}{5}\sqrt{(5x^2+8)^3} + k_1.
\end{align*}

Para a segunda integral, pode-se fazer a substituição $w = \sen(x)$, $\diff w = \cos(x) \diff x$:
\begin{align*}
\int \cotg(x) \diff x
&= \int \frac{\cos(x)}{\sen(x)} \diff x
= \int \frac{1}{\sen(x)} \cos(x)\diff x
= \int \frac{1}{w} \diff w\\
&= \ln|w| + k_2
= \ln|\sen(x)| + k_2.
\end{align*}
Portanto, $\int 3x \sqrt{5x^2+8} -\cotg(x) \diff x = \frac{1}{5}\sqrt{(5x^2+8)^3} - \ln|\sen(x)| + c$.

\item Decompondo a função em frações parciais, e substituindo $u=\frac{x}{3}$, $\diff x = 3 \diff u$, $v=\frac{x}{2}$ e $\diff x = 2 \diff v$, resulta que:
\begin{align*}
  \int \frac{25}{(x^2+9)(x^2+4)} \diff x
& = \int \frac{-5}{x^2+9} + \frac{5}{x^2+4} \diff x
  = -\frac{5}{9} \int \frac{1}{\left(\frac{x}{3}\right)^2+1} \diff x
   + \frac{5}{4} \int\frac{1}{\left(\frac{x}{2}\right)^2+1} \diff x\\
& = -\frac{5}{9} \int \frac{1}{u^2+1} 3\diff u
  + \frac{5}{4} \int\frac{1}{v^2+1} 2\diff v\\
& = -\frac{5}{3} \arctg(u) + \frac{5}{2} \arctg(v)+c
  = -\frac{5}{3} \arctg\left(\frac{x}{3}\right) + \frac{5}{2} \arctg\left(\frac{x}{2}\right)+c.
\end{align*}

\item Usando propriedades da função exponencial e substituindo $z=e^x$ e $\diff z = e^x \diff x$, tem-se:
\begin{align*}
  \int \frac{e^{(x+3)}}{1+e^{(2x)}} \diff x
&= \int \frac{e^3e^x}{1+(e^x)^2} \diff x
= e^3 \int \frac{1}{1+(e^x)^2} e^x \diff x
= e^3 \int \frac{1}{1+z^2} \diff z\\
&= e^3 \arctg(z)+c
= e^3 \arctg(e^x)+c.
\end{align*}

\item Decompondo a função em frações parciais, resulta que:
\begin{align*}
  \int \frac{9 x^2 - 14 x + 2}{x^3-2x^2+x} \diff x
& = \int \frac{9 x^2 - 14 x + 2}{x(x-1)^2} \diff x
  = \int \frac{2}{x} + \frac{7}{x-1} + \frac{25}{(x-1)^2} \diff x\\
& = 2\int \frac{1}{x} \diff x
  + 7\int \frac{1}{x-1} \diff x
 + 25\int \frac{1}{(x-1)^2} \diff x\\
& = 2\ln|x| + 7\ln|x-1| - \frac{25}{x-1} + c\\
& = \ln|x^2(x-1)^7| - \frac{25}{x-1} + c.
\end{align*}

\end{enumerate}
\end{enumerate}
\end{document}
